% Cryptographically Keyed GDSS -- AI-generated preprint
% Compile: cd paper && pdflatex -interaction=nonstopmode kgdss_paper.tex && pdflatex kgdss_paper.tex
\documentclass[10pt,a4paper,twocolumn]{article}
\usepackage[utf8]{inputenc}
\usepackage[T1]{fontenc}
\usepackage{mathptmx}
\usepackage{amsmath,amssymb}
\usepackage{graphicx}
\usepackage{booktabs}
\usepackage{xcolor}
\usepackage{titlesec}
\usepackage{fancyhdr}
\usepackage{url}
\usepackage{hyperref}
\usepackage{caption}
\usepackage{float}
\usepackage{dblfloatfix}
\usepackage[margin=2.5cm]{geometry}

\definecolor{mdpiblue}{HTML}{004E8F}
\titleformat{\section}{\normalfont\bfseries\color{mdpiblue}\large}{\thesection.}{0.4em}{}
\titleformat{\subsection}{\normalfont\bfseries\normalsize}{\thesubsection}{0.4em}{}
\captionsetup{font=small,labelfont=bf}

\pagestyle{fancy}
\fancyhf{}
\fancyhead[L]{\small Sensors 2025, xx(xx), xxxx}
\fancyhead[R]{\small doi:10.3390/xxxx (placeholder)}
\fancyfoot[C]{\thepage}
\renewcommand{\headrulewidth}{0.4pt}
\hypersetup{colorlinks=true,linkcolor=mdpiblue,urlcolor=mdpiblue}

\newcommand{\fig}[1]{figures/#1}
\linespread{1.12}

\begin{document}

\onecolumn
\noindent\setlength{\fboxrule}{1.5pt}\setlength{\fboxsep}{8pt}%
\fcolorbox{black}{yellow!85!white}{%
\parbox{\dimexpr\textwidth-2\fboxsep-2\fboxrule\relax}{%
\textbf{DISCLAIMER:} This paper was generated in its entirety by an artificial
intelligence system (Claude, by Anthropic) from source code, test results,
and design documentation. It has not been reviewed by a professional
cryptographer, signals intelligence specialist, or academic peer reviewer.
The theoretical and security claims presented here are based on logical
reasoning from published sources and should be treated with appropriate
scepticism. This paper is provided for research and discussion purposes
only. Independent expert review is strongly recommended before any
practical application.}%
}
\vspace{14pt}
\begin{center}
{\Large\bfseries Cryptographically Keyed Gaussian-Distributed Spread-Spectrum\\
for Enhanced Covert Communications: Design, Implementation,\\
and Simulated Performance in ITU Channel Models}\\[10pt]
{\normalsize\textbf{Status:} Preprint --- AI-generated, unreviewed}\\[6pt]
{\small\textit{Journal style only: Sensors (MDPI), Section: Communications --- not a submission.}}\\[4pt]
{\small ISSN 1424-8220 (style placeholder)}\\[8pt]
{\small\textbf{Keywords:} covert communications; spread-spectrum; low probability of detection;
cryptographic key derivation; ECDH; ChaCha20; SOQPSK; GNU Radio; LDPC;
Brainpool curves; Shamir secret sharing; hardware security module;
software-defined radio}
\end{center}
\vspace{10pt}
\begin{abstract}
\noindent
Gaussian-Distributed Spread-Spectrum (GDSS) achieves noise-like statistics for low probability of detection (LPD), yet the original design leaves a practical traffic-analysis surface: synchronisation bursts built from repeatable pseudo-noise (PN) and predictable timing can correlate across sessions. Masking drawn from thermal noise is statistically Gaussian but not cryptographically secret against a determined adversary with long captures. This preprint documents \emph{Cryptographically Keyed GDSS} (GR-K-GDSS), implemented as the GNU Radio out-of-tree module \texttt{gr-k-gdss}, which replaces hardware-noise masking with ChaCha20-driven Box--Muller Gaussian masks, derives independent subkeys via HKDF-SHA256 from a BrainpoolP256r1 ECDH shared secret, and randomises sync PN and timing per session. Session nonces concatenate a 32-bit session identifier and 64-bit transmit sequence in big-endian form. The reference transmit chain (example flowgraph) chains microphone input, Codec2 vocoder, ChaCha20-Poly1305 payload protection (via \texttt{gr-linux-crypto}), SOQPSK modulation (\texttt{gr-qradiolink}), and the keyed spreader. The despreader implements acquisition, tracking, and locked states with coarse-fine code search, early-prompt-late timing, adaptive correlation thresholds, SNR estimation, and phase-based frequency error tracking; keys arrive on a GNU Radio PMT \texttt{set\_key} port. Simulated IQ tests in the repository report cross-session standard-GDSS sync correlation 1.0000 versus 0.1028 for keyed bursts ($\approx$9.7$\times$ reduction). Unit tests: 37 passed (1 skipped) at time of generation. IQ statistical checks: 29/29 passed on generated files. Section~7 reports statistical bit-error simulations (NumPy/SciPy Monte Carlo for standard and keyed GDSS, DSSS analytical reference, simplified VHF/HF channels); they are not over-the-air measurements. LDPC overlays use ideal SNR shifts ($\sim$5~dB) consistent with prior GDSS/LDPC studies rather than bit-true decoder output. Limitations include absence of formal indistinguishability proofs for ChaCha20+Box--Muller versus thermal noise, lack of cryptographic forward secrecy under static long-term ECDH keys unless augmented, and quantum vulnerability of BrainpoolP256r1 to Shor's algorithm. Operational guidance for RTL-SDR DC offset and IQ imbalance appears in project \texttt{USAGE.md}. All claims are contingent on expert review.
\end{abstract}
\vspace{16pt}

\twocolumn

\section{Introduction}
\subsection{Background}
Covert or LPD communications seek to hide the presence of transmissions from passive observers, not only ciphertext from eavesdroppers \cite{bash2015,bash2013,turner1991}.

\subsection{Standard DSSS detectability}
Direct-sequence spread spectrum reduces power spectral density but retains cyclostationary structure exploitable by cyclic feature detectors \cite{gardner1992}. Shakeel et al.\ summarise DSSS limitations and propose GDSS \cite{shakeel2023}.

\subsection{Standard GDSS and residual issues}
GDSS masks each chip with Gaussian amplitudes so higher-order moments match thermal noise \cite{shakeel2023}. Standard GDSS draws masks from transmitter thermal noise; the algorithm is public and masking is statistically constrained but not key-secret. Sync bursts in the reference design can reuse PN and timing, enabling cross-session correlation (traffic analysis).

\subsection{Motivation for cryptographic keying}
Keyed masking makes the chip-wise Gaussian amplitudes unpredictable without the session key, while preserving first-order noise-like appearance. Session-specific PN and timing for sync bursts reduces repeatable energy across sessions \cite{repo_gdss}.

\subsection{Contributions}
We (1) specify HKDF-separated subkeys for payload AEAD, GDSS masking, sync PN, and sync timing; (2) document the \texttt{gr-k-gdss} spreader/despreader and Python helpers; (3) summarise repository tests and statistical BER simulations; (4) compare detection and jamming scores tabulated in the project README. Section~2 reviews standard GDSS; Section~3 specifies keyed GDSS; Section~4 implementation; Section~5 security discussion; Section~6 tests; Section~7 statistical BER simulations; Sections~8--9 discuss limitations and conclude.

\section{Standard GDSS --- Review}
\subsection{DSSS}
Symbols are multiplied by a spreading sequence; processing gain scales with spreading factor $N$ \cite{shakeel2023}.

\subsection{GDSS masking}
For complex symbol $S$ and independent Gaussian $U,V$,
\begin{equation}
I = \mathrm{Re}(S)\,|U|,\quad Q = \mathrm{Im}(S)\,|V|.
\label{eq:gdssmask}
\end{equation}
Quadrant information remains in the sign of $\mathrm{Re}(S)$ and $\mathrm{Im}(S)$ while amplitudes follow folded-Gaussian statistics \cite{shakeel2023}.

\subsection{Gaussianity (Table~\ref{tab:moments})}
Shakeel et al.\ match even moments through order 20 between GDSS and theoretical Gaussian; DSSS deviates (Table~\ref{tab:moments}).

\begin{table*}[t]
\centering
\caption{Moments of noise-free signals ($N{=}256$), after Shakeel et al.\ \cite{shakeel2023}.}
\label{tab:moments}
\small
\begin{tabular}{@{}lrrr@{}}
\toprule
Moment & SGD (theory) & GDSS & DSSS \\
\midrule
2nd & 1 & 1 & 1 \\
4th & 3 & 3 & 1 \\
6th & 15 & 15 & 1 \\
8th & 105 & 105 & 1 \\
10th & 945 & 947 & 1 \\
12th & 10{,}395 & 10{,}428 & 1 \\
14th & 135{,}135 & 135{,}524 & 1 \\
16th & 2{,}027{,}025 & 2{,}025{,}880 & 1 \\
18th & 34{,}459{,}425 & 34{,}102{,}797 & 1 \\
20th & 654{,}729{,}075 & 634{,}790{,}833 & 1 \\
\bottomrule
\end{tabular}
\end{table*}

\subsection{Detectability}
Moment tests, modulation stripping, and cyclostationary analysis favour GDSS over DSSS at low SNR \cite{shakeel2023}.

\subsection{Sync burst vulnerability}
Fixed PN and timing let a long-term observer align bursts across sessions (correlation $\approx 1$ in software simulation).

\section{Proposed Keyed GDSS Scheme}
\subsection{Goals}
Maintain noise-like statistics while making masks, sync PN, and timing cryptographically session-specific.

\subsection{Key derivation}
BrainpoolP256r1 ECDH yields a shared secret; HKDF-SHA256 \cite{rfc5869} expands four 32-byte keys with distinct \texttt{info} strings (Table~\ref{tab:hkdf}).

\begin{table}[H]
\centering
\caption{HKDF subkeys (\texttt{session\_key\_derivation.py}).}
\label{tab:hkdf}
\small
\begin{tabular}{@{}llp{4.2cm}@{}}
\toprule
Name & Info string & Purpose \\
\midrule
\texttt{payload\_enc} & \texttt{payload-chacha20poly1305-v1} & Payload AEAD \\
\texttt{gdss\_masking} & \texttt{gdss-chacha20-masking-v1} & Chip masks \\
\texttt{sync\_pn} & \texttt{sync-dsss-pn-sequence-v1} & Sync PN \\
\texttt{sync\_timing} & \texttt{sync-burst-timing-offset-v1} & Timing \\
\bottomrule
\end{tabular}
\end{table}

\begin{figure*}[t]
\centering
\includegraphics[width=0.92\textwidth]{\fig{fig2_key_derivation.png}}
\caption{Key derivation hierarchy (conceptual).}
\end{figure*}

\subsection{Cryptographic masking}
Sixteen ChaCha20 bytes per chip provide four 32-bit uniforms; Box--Muller yields two Gaussian masks:
\begin{equation}
z=\sqrt{-2\ln u_1}\cos(2\pi u_2).
\end{equation}
Minimum magnitude clamp $\mathrm{MIN\_MASK}=10^{-4}$ avoids division instability. Transmitted chip (conceptually) follows \eqref{eq:gdssmask} with keyed Gaussians; receiver divides by the same masks then sums:
\begin{equation}
\hat{z}=\frac{1}{N}\sum_{i=1}^{N}\frac{r_i}{m_i}.
\end{equation}
Implementation matches \texttt{kgdss\_spreader\_cc} / \texttt{kgdss\_despreader\_cc} and libsodium ChaCha20 IETF interface \cite{rfc8439,bernstein2008}.

\subsection{Session-unique sync}
\texttt{derive\_sync\_pn\_sequence}: HMAC-SHA256(\texttt{sync\_pn}, \texttt{sync-pn-v1}$\Vert$\texttt{session\_id}) derives material expanded by ChaCha20 to $\pm1$ chips \cite{repo_gdss}.\\
\texttt{derive\_sync\_schedule}: derives per-epoch offsets in $[-w,w]$ ms via ChaCha20-indexed uniforms \cite{repo_gdss}.\\
\texttt{gaussian\_envelope}: rise/fall $\exp(-x^2/2)$ on flanks, \texttt{rise\_fraction}=0.1.

\begin{figure*}[t]
\centering
\includegraphics[width=0.88\textwidth]{\fig{fig4_sync_burst.png}}
\caption{Conceptual standard vs keyed sync placement.}
\end{figure*}

\subsection{Signal chain}
\begin{figure*}[t]
\centering
\includegraphics[width=0.95\textwidth]{\fig{fig1_block_diagram.png}}
\caption{Transmit stack and keying path (after \texttt{tx\_example\_kgdss.grc}).}
\end{figure*}

\section{Implementation}
\subsection{Modules}
\texttt{gr-k-gdss}: C++ spreader/despreader/key injector; Python bindings.\\
\texttt{gr-qradiolink}: SOQPSK, legacy GDSS, LDPC.\\
\texttt{gr-linux-crypto}: ECDH, AEAD, ECIES, Shamir, keyring, Nitrokey hooks \cite{repo_crypto}.

\subsection{Despreader}
States \texttt{ACQUISITION}, \texttt{TRACKING}, \texttt{LOCKED}; \texttt{COARSE\_SEARCH\_BINS}=32; \texttt{LOCK\_THRESHOLD}=10; \texttt{ADAPTIVE\_THRESHOLD\_MIN}=0.2; early-prompt-late timing error scaled by 0.1; SNR IIR $\alpha=0.05$; frequency error from correlation phase diffs with $\alpha=0.1$; mutexes \texttt{d\_mutex}, \texttt{d\_key\_mutex}; PMT \texttt{set\_key} \cite{repo_gdss}.

\begin{figure}[H]
\centering
\includegraphics[width=\linewidth]{\fig{fig3_state_machine.png}}
\caption{Despreader state machine (simplified).}
\end{figure}

\subsection{Python helpers}
\texttt{derive\_session\_keys}, \texttt{store\_session\_keys}, \texttt{gdss\_nonce}, \texttt{payload\_nonce}, \texttt{get\_shared\_secret\_from\_gnupg} \cite{repo_gdss}.

\subsection{Cryptographic substrate}
OpenSSL/cryptography; Brainpool ECDH/ECDSA; multi-recipient ECIES; Shamir; Nitrokey interface; Linux keyring storage \cite{repo_crypto,bsi2018}.

\subsection{SOQPSK and LDPC}
Operating modes and LDPC rates/block lengths per \texttt{gr-qradiolink} documentation \cite{repo_qrl,chung2001}.

\section{Security Analysis}
\subsection{Detection resistance}
Table~\ref{tab:det} reproduces README 0--10 ratings. Largest delta: sync burst and traffic analysis \cite{repo_gdss}.

\begin{table}[H]
\centering
\caption{Detection resistance (README).}
\label{tab:det}
\small
\begin{tabular}{@{}lccc@{}}
\toprule
Aspect & Std & Keyed & $\Delta$ \\
\midrule
Passive detection & 7 & 7 & 0 \\
Cyclostationary & 8 & 8 & 0 \\
Moments & 8 & 8 & 0 \\
Mod.\ stripping & 8 & 8 & 0 \\
Sync burst & 4 & 8 & +4 \\
Traffic analysis & 5 & 8 & +3 \\
Noise floor & 5 & 6 & +1 \\
Direction finding & 5 & 5 & 0 \\
\midrule
\textbf{Overall} & \textbf{6} & \textbf{7.5} & \textbf{+1.5} \\
\bottomrule
\end{tabular}
\end{table}

\begin{figure}[H]
\centering
\includegraphics[width=\linewidth]{\fig{fig11_radar.png}}
\caption{Spider chart (normalised).}
\end{figure}

\subsection{Content recovery and jamming}
See README tables (decryption/masking strip; broadband and sync-targeted jamming).

\subsection{Keys and forward secrecy}
Static long-term ECDH keys imply no formal forward secrecy; Nitrokey PIN, brick policy, and physical destruction mitigate compromise \cite{repo_gdss}. Ephemeral ECDH authenticated by long-term keys is recommended.

\subsection{Quantum}
BrainpoolP256r1 is Shor-vulnerable; symmetric layers degrade under Grover; PQ KEM hook raises \texttt{NotImplementedError} until enabled \cite{repo_crypto}.

\section{Verification and Tests}
\subsection{Unit tests}
At generation time: \textbf{37 passed, 1 skipped} in \texttt{tests/} (\texttt{test\_cross\_layer}, \texttt{test\_t1\_spreader\_despreader}, \texttt{test\_t2\_sync\_burst}, \texttt{test\_t3\_key\_derivation}). README snapshots previously cited 30 tests; the suite has since expanded.

\subsection{IQ analysis}
\textbf{29/29} checks passed on generated files (\texttt{docs/TEST\_RESULTS.md}): Gaussian baseline and keyed transmission moments, wrong-key isolation, nonce reuse warning, standard GDSS statistics, KL vs keyed, keyed cross-session peak $<0.15$.

\begin{figure}[H]
\centering
\includegraphics[width=\linewidth]{\fig{fig5_cross_session_bar.png}}
\caption{Cross-session sync correlation (repository generator output).}
\label{fig:cross}
\end{figure}

\subsection{gr-linux-crypto}
NIST CAVP vectors (AES-GCM, ChaCha20-Poly1305), Wycheproof ECDH/ECDSA on Brainpool curves, BSI TR-03111 (20/20), ECTester (24 passed, 1 skipped), 805M+ LibFuzzer executions without crash, CBMC checks, dudect timing, microsecond-level AEAD latency \cite{repo_crypto_results}.

\subsection{gr-qradiolink}
C++ and MMDVM tests and fuzzing summarised in project documentation \cite{repo_qrl}.

\begin{table}[H]
\centering
\caption{Illustrative crypto validation snapshot (see gr-linux-crypto TEST\_RESULTS).}
\small
\begin{tabular}{@{}lcc@{}}
\toprule
Suite & Result \\
\midrule
NIST AES-128-GCM & 4/4 \\
NIST AES-256-GCM & 4/4 \\
RFC8439 ChaCha20-Poly1305 & 3/3 \\
Wycheproof Brainpool ECDH & 2534+ vectors pass \\
BSI TR-03111 & 20/20 \\
\bottomrule
\end{tabular}
\end{table}

\begin{table}[H]
\centering
\caption{Decryption/content recovery resistance (README).}
\scriptsize
\begin{tabular}{@{}lccc@{}}
\toprule
Aspect & Std & Keyed & $\Delta$ \\
\midrule
Payload decryption & 9 & 9 & 0 \\
Masking strip attack & 4 & 9 & +5 \\
Statistical masking recovery & 5 & 9 & +4 \\
Forward secrecy (rated) & 7 & 9 & +2 \\
Key compromise blast radius & 6 & 8 & +2 \\
\midrule
\textbf{Overall} & \textbf{6} & \textbf{9} & \textbf{+3} \\
\bottomrule
\end{tabular}
\end{table}

\begin{table}[H]
\centering
\caption{Jamming resistance (README).}
\scriptsize
\begin{tabular}{@{}lccc@{}}
\toprule
Aspect & Std & Keyed & $\Delta$ \\
\midrule
Broadband noise jamming & 6 & 6 & 0 \\
Targeted carrier jamming & 8 & 8 & 0 \\
Protocol-aware jamming & 5 & 7 & +2 \\
Sync burst jamming & 4 & 7 & +3 \\
Replay jamming & 5 & 7 & +2 \\
\midrule
\textbf{Overall} & \textbf{5.5} & \textbf{7} & \textbf{+1.5} \\
\bottomrule
\end{tabular}
\end{table}

\begin{table}[H]
\centering
\caption{Key compromise scenarios (README).}
\scriptsize
\begin{tabular}{@{}p{3.2cm}p{3.5cm}p{3.8cm}@{}}
\toprule
Scenario & Protection & Outcome \\
\midrule
Device seized locked & PIN + auto-brick & Key inaccessible \\
Device active/unlocked & Keyring in use & Session exposure risk \\
Compelled PIN + disposal option & Physical destruction & Key destroyed \\
Destroyed before seizure & None left & Past sessions unrecoverable \\
\bottomrule
\end{tabular}
\end{table}

\begin{table*}[t]
\centering
\caption{SOQPSK operating parameters (representative; see gr-qradiolink README).}
\small
\begin{tabular}{@{}lllll@{}}
\toprule
Mode & Application & Data rate & BW / spacing & Notes \\
\midrule
VHF Mode 1 & Land mobile & 14.4 kbps & $\sim$10 kHz / 12.5 kHz ch. & Single carrier \\
VHF Mode 2 & Multipath & $3\times$4.8 kbps & 4 kHz carrier spacing & Multicarrier \\
HF standard & Narrow HF & 4 kbps/carrier & 2.7 kHz & STANAG-style bands \\
HF 10m & Wide HF & 9 kbps/carrier & 6 kHz & Higher-rate option \\
\bottomrule
\end{tabular}
\end{table*}

\begin{table*}[t]
\centering
\caption{gr-linux-crypto cross-validation (summary from TEST\_RESULTS.md).}
\scriptsize
\begin{tabular}{@{}llrrr@{}}
\toprule
Suite & Algorithm / scope & Vectors & Passed & Rate \\
\midrule
NIST CAVP & AES-128-GCM & 4 & 4 & 100\% \\
NIST CAVP & AES-256-GCM & 4 & 4 & 100\% \\
NIST CAVP & RFC 8439 ChaCha20-Poly1305 & 3 & 3 & 100\% \\
Wycheproof & ECDH Brainpool P-256r1 & 2534+ & 2534+ & 100\% \\
Wycheproof & ECDH Brainpool P-384r1 & (set) & all & 100\% \\
Wycheproof & ECDH Brainpool P-512r1 & (set) & all & 100\% \\
Wycheproof & ECDSA per Brainpool curve & 475+ each & 475+ & 100\% \\
\bottomrule
\end{tabular}
\end{table*}

\begin{table}[H]
\centering
\caption{Representative AEAD latency (gr-linux-crypto benchmarks, $\mu$s).}
\small
\begin{tabular}{@{}lrrrr@{}}
\toprule
Algorithm & Mean & p50 & p95 & Thr. \\
\midrule
AES-128-GCM & 8.8 & -- & -- & $<$100 \\
AES-256-GCM & 9.3 & -- & -- & $<$100 \\
ChaCha20-Poly1305 & 11.5 & -- & -- & $<$100 \\
\bottomrule
\end{tabular}
\end{table}

\begin{table*}[t]
\centering
\caption{IQ file analysis outcomes (\texttt{analyse\_iq\_files.py}, 29 checks PASS).}
\scriptsize
\begin{tabular}{@{}ll@{}}
\toprule
File / comparison & Tests (all PASS) \\
\midrule
\texttt{01\_gaussian\_noise\_baseline.cf32} & Mean I/Q, variance symmetry, kurtosis I/Q, skewness I/Q, autocorr. \\
\texttt{03\_keyed\_gdss\_transmission.cf32} & Same eight statistics \\
\texttt{04\_keyed\_gdss\_despread\_correct\_key.cf32} & Round-trip correlation \\
\texttt{05\_keyed\_gdss\_despread\_wrong\_key.cf32} & Key isolation \\
\texttt{07\_nonce\_reuse} & XOR correlation / reuse warning \\
\texttt{09\_standard\_gdss\_transmission.cf32} & Eight noise-like tests \\
\texttt{09\_vs\_03} & KL divergence (I) \\
\texttt{13\_keyed\_gdss\_crosscorr} & Cross-session peak $<0.15$ \\
\bottomrule
\end{tabular}
\end{table*}

\begin{table*}[t]
\centering
\caption{gr-k-gdss pytest modules (snapshot: 37 passed, 1 skipped).}
\scriptsize
\begin{tabular}{@{}ll@{}}
\toprule
Module / test & Status \\
\midrule
\texttt{test\_cross\_layer.py::test\_full\_stack\_round\_trip} & PASSED \\
\texttt{test\_t1\_spreader\_despreader.py} (14 tests: round\_trip, keystream, key/nonce sensitivity, invalid sizes, Gaussian, masks, boundary, set\_key) & PASSED \\
\texttt{test\_t2\_sync\_burst.py} (12 tests: PN, timing, envelope, masking) & PASSED \\
\texttt{test\_t3\_key\_derivation.py} (10 tests: HKDF, nonces, keyring) & PASSED \\
Optional keyring / environment & 1 skipped \\
\bottomrule
\end{tabular}
\end{table*}

\section{BER Performance (Statistical Simulation)}
\subsection{Simulation methodology}
Following Shakeel et al., the DSSS reference is
\begin{equation}
P_b = \tfrac12\operatorname{erfc}\!\sqrt{N E_s/(2N_0)}.
\label{eq:dsss}
\end{equation}
Curves in Figs.~\ref{fig:ber7}--\ref{fig:ber10} are produced by \texttt{paper/ber\_simulation.py} (NumPy/SciPy), aligned with the GNU Radio stack conceptually but without invoking the SDR runtime. Spreading factors $N\in\{64,128,256\}$ match the implementation. The SNR axis is $E_b/N_0$ in dB from $-20$ to $+5$ in 1~dB steps. Unless overridden by the environment variable \texttt{BER\_MC\_NUM\_BITS}, each plotted point uses $10^6$ simulated bits per receiver condition.

Standard GDSS uses chips $x_i=s|G_i|$ with $G_i\sim\mathcal{N}(0,1)$ and decision $\operatorname{sign}(\sum_i r_i)$. Keyed GDSS uses Box--Muller masks $m_i$ (clamped to minimum magnitude $10^{-4}$ as in the C++ block) and $\operatorname{sign}(\operatorname{mean}_i(r_i/m_i))$. AWGN variance per chip follows the repetition-combine DSSS mapping so that the DSSS curve matches \eqref{eq:dsss}.

VHF land mobile (ITU-R P.1406--inspired): Rayleigh amplitude per symbol, Doppler-induced phase walk across chips at 50~Hz (pedestrian) and 200~Hz (vehicular); SOQPSK Mode~2 is approximated by a $+0.6$~dB effective $E_b/N_0$ offset relative to Mode~1 (multicarrier diversity placeholder). LDPC rate $1/2$ overlays use ideal waterfall shifts ($\approx$5~dB on keyed curves) consistent with \cite{shakeel2023} and gr-qradiolink usage \cite{repo_qrl,chung2001}.

HF (STANAG~4539--inspired): tapped-delay-line profiles labelled AWGN, Good, Poor, and Disturbed convolve chip sequences before AWGN; comparisons use uncoded standard GDSS versus LDPC-shifted keyed GDSS as in Fig.~\ref{fig:ber9}.

Semi-analytical keyed extension (conceptual):
\begin{equation}
P_b \approx \int_{-\infty}^{0} p_{\hat{Z}}(z)\,dz,
\end{equation}
with $\hat{Z}$ the $N$-fold combination after per-chip division by Gaussian masks; closed forms are not used here---Monte Carlo estimates $P_b$ directly.

HKDF expand (RFC~5869) satisfies
\begin{equation}
T(i)=\mathrm{HMAC\text{-}SHA256}(\mathrm{PRK},\,T(i{-}1)\,\Vert\,\mathrm{info}\,\Vert\,i),\quad \mathrm{OKM}=T(1)\,\Vert\,T(2)\,\Vert\cdots
\end{equation}
(first $L$ octets). GDSS ChaCha20 nonce:
\begin{equation}
\mathrm{nonce}=\mathrm{BE}_{4}(\mathrm{session\_id})\,\Vert\,\mathrm{BE}_{8}(\mathrm{tx\_seq}).
\end{equation}
Envelope flank samples use $x\in[-3,0]$ linearly spaced with $\exp(-x^2/2)$ normalised; adaptive lock threshold $\theta_{\mathrm{adaptive}}=\max(0.2,\,\theta\cdot C_{\mathrm{avg}}/C_{\mathrm{peak}})$ with $C_{\mathrm{avg}}$ IIR-smoothed correlation magnitude.

\subsection{AWGN baseline}
Fig.~\ref{fig:ber7} contrasts \eqref{eq:dsss} with Monte Carlo standard and keyed GDSS. At mid-BER, keyed GDSS typically sits $\sim$2--2.5~dB to the right of standard GDSS, in line with Shakeel et al.\ Figure~14-style mask-processing loss.

\subsection{VHF land mobile channel}
Fig.~\ref{fig:ber8} fixes $N{=}256$ and shows Mode~1 vs.\ Mode~2 with and without ideal LDPC overlays, for 50~Hz and 200~Hz max Doppler.

\subsection{HF military bands}
Fig.~\ref{fig:ber9} uses $N{=}256$ across four stylised HF profiles.

\subsection{LDPC coding gain}
Fig.~\ref{fig:ber10} compares uncoded keyed AWGN ($N{=}256$) to ideal rate-$1/2$ gains of 4.8~dB (block 576) and 5.2~dB (block 1152), consistent with Shakeel et al.\ Figure~15 ordering.

\begin{figure*}[t]
\centering
\includegraphics[width=0.92\textwidth]{\fig{fig7_awgn_ber.png}}
\caption{AWGN: DSSS \eqref{eq:dsss} vs Monte Carlo standard and keyed GDSS ($N=64,128,256$).}
\label{fig:ber7}
\end{figure*}

\begin{figure*}[t]
\centering
\includegraphics[width=0.92\textwidth]{\fig{fig8_vhf_ber.png}}
\caption{VHF-style Rayleigh fading with Doppler phase walk; $N{=}256$, SOQPSK Mode~1 vs.\ Mode~2, uncoded and LDPC-shifted keyed GDSS.}
\label{fig:ber8}
\end{figure*}

\begin{figure*}[t]
\centering
\includegraphics[width=0.92\textwidth]{\fig{fig9_hf_ber.png}}
\caption{HF-style tapped-delay profiles (AWGN, Good, Poor, Disturbed): uncoded standard GDSS vs LDPC-coded keyed GDSS ($N{=}256$).}
\label{fig:ber9}
\end{figure*}

\begin{figure*}[t]
\centering
\includegraphics[width=0.88\textwidth]{\fig{fig10_ldpc.png}}
\caption{Keyed GDSS on AWGN: uncoded vs ideal LDPC rate $1/2$ with 576- and 1152-bit block gain anchors ($N{=}256$).}
\label{fig:ber10}
\end{figure*}

\begin{figure*}[t]
\centering
\includegraphics[width=0.95\textwidth]{\fig{fig6_histograms.png}}
\caption{Histogram comparison (synthetic stand-in for IQ report files).}
\end{figure*}

\section{Discussion}
\subsection{Comparison to standard GDSS}
Keyed GDSS closes the masking-strip gap and strengthens sync obscurity; PAPR remains high as in standard GDSS (Gaussian amplitudes). Operational security (short transmissions, mobility) complements physics limits (direction finding).

\subsection{Literature positioning}
Featureless-waveform concepts date to early LPI research \cite{turner1991}; cyclostationary countermeasures \cite{gardner1992} motivate the GDSS masking approach \cite{shakeel2023}. Chaos-based and artificial-noise schemes \cite{kaddoum2016,soltani2018} pursue alternative statistical hides. GR-K-GDSS is closest in spirit to GDSS but swaps thermal-noise draws for ChaCha20-Box--Muller outputs and adds sessionised sync.

\subsection{PAPR and hardware}
Gaussian chip amplitudes imply high peak-to-average power ratio relative to constant-envelope modulations. Power amplifiers may clip unless backed off, reintroducing distortion features; this mirrors standard GDSS hardware constraints \cite{shakeel2023}.

\subsection{Open problems}
Formal indistinguishability games for ChaCha20-derived masks, complete OTA validation with Nitrokey-backed keys, and integration of ephemeral ECDH for forward secrecy remain outstanding \cite{repo_gdss}.

\section{Conclusions}
GR-K-GDSS integrates HKDF-separated keys, ChaCha20 Box--Muller masking, and sessionised sync helpers into GNU Radio. Repository tests support functional correctness and large cross-session correlation reduction in simulation. Future work: ephemeral ECDH, PQ KEM integration, OTA validation with hardware tokens, formal statistical analysis of ChaCha20-Gaussian outputs.

\clearpage
\onecolumn
\section*{Appendix A: Literature context for LPD and covert waveforms}
Chaos-based and artificial-noise assisted covert schemes \cite{kaddoum2016,soltani2018} illustrate alternative routes to hiding power or structure; GDSS instead matches thermal statistics directly. Surveys such as \cite{makhdoom2022} catalogue detection metrics relevant when evaluating keyed versus standard masking. Fundamental limits from information-theoretic covert communication \cite{bash2013,bash2015} bound how much data can be sent covertly under warden models; this engineering design does not claim to meet those bounds but inherits spread-spectrum processing gain.

\section*{Appendix B: Operational notes from repository documentation}
SDR hardware often exhibits a strong DC component and IQ imbalance. The GR-K-GDSS blocks do not correct these; GNU Radio users should apply offset tuning, DC blocking, and IQ balance blocks ahead of the despreader, as summarised in \texttt{docs/USAGE.md}. Plot scripts subtract means for visualisation only.

\section*{Appendix C: Repository map}
\texttt{lib/} C++ spreader and despreader; \texttt{python/} session and sync helpers; \texttt{grc/} block YAML; \texttt{tests/} IQ generators, \texttt{analyse\_iq\_files.py}, \texttt{plot\_iq\_comparison.py}, \texttt{plot\_spectrum\_snapshots.py}; \texttt{examples/tx\_example\_kgdss.grc} reference flowgraph; \texttt{docs/} USAGE, TESTING, TEST\_RESULTS, GLOSSARY.

\section*{Appendix D: Reproducibility}
Clone \cite{repo_gdss}; install GNU Radio 3.10+, libsodium, gr-linux-crypto, gr-qradiolink as dependencies; build with CMake. Generate IQ fixtures: \texttt{python3 tests/generate\_iq\_test\_files.py}. Analyse: \texttt{python3 tests/analyse\_iq\_files.py}. BER Monte Carlo data: \texttt{cd paper \&\& python3 ber\_simulation.py} (optional: \texttt{BER\_MC\_NUM\_BITS=50000} for a faster dry run). Figures: \texttt{python3 gen\_figures.py} (loads \texttt{paper/figures/ber\_mc\_results.npz} for Figs.~\ref{fig:ber7}--\ref{fig:ber10}; parametric fallbacks apply if the file is absent). PDF: \texttt{pdflatex kgdss\_paper.tex} (twice).

\vspace{1em}
\noindent\textbf{Limitation reminder.} Section~7 uses statistical baseband models (not over-the-air captures). VHF and HF channels are simplified Rayleigh and tapped-delay abstractions; they are not full ITU-R P.1406 or STANAG~4539 reference simulators. LDPC curves apply ideal SNR shifts to the keyed uncoded Monte Carlo BER, not bit-true decoder output. Independent replication with waveform-accurate GNU Radio flowgraphs is required before quoting absolute link budgets.

\vspace{0.8em}
\noindent\textbf{Channel simulation disclaimer.} ITU-R P.1406 models terrestrial land-mobile propagation with height, clutter, and frequency-dependent path loss \cite{itu1406}. A faithful implementation couples geometry databases, antenna patterns, and time-varying mobile routes. Section~7.3 uses Rayleigh fading with chip-wise Doppler phase rotation at labelled max Doppler (50~Hz and 200~Hz), not the standard's full scenario classes. STANAG~4539 HF models specify tapped-delay-line profiles for poor, good, and disturbed conditions \cite{stanag4539}; Fig.~\ref{fig:ber9} uses stylised normalised taps and real-baseband BPSK-on-$I$ chip ISI before AWGN, not a full STANAG modem chain.

\vspace{0.5em}
\noindent\textbf{Comparison to Shakeel et al.\ Figure~14--15.} The original GDSS paper contrasts DSSS, GDSS, and LDPC-coded GDSS using full waveform simulation \cite{shakeel2023}. This preprint adds Monte Carlo chip-level keyed GDSS and ideal LDPC overlays; aligning every decibel with \cite{shakeel2023} still requires bit-matching SOQPSK pulse shapes, resampler group delay, and LDPC interleaving as built in gr-qradiolink \cite{repo_qrl}.

\vspace{0.5em}
\noindent\textbf{Side-channel and fuzzing scope.} gr-linux-crypto aggregates NIST CAVP downloads, Wycheproof JSON vectors, BSI TR-03111 structural checks, ECTester compatibility suites, LibFuzzer targets with ASan/UBSan, CBMC slice verification, and dudect-style timing experiments \cite{repo_crypto_results,reparaz2017}. These results bound defects in OpenSSL/cryptography call paths wrapped by the module. They do not cover the fixed-point behaviour of VOLK kernels inside GNU Radio, analogue front-end leakage, or electromagnetic emanations from SDR hardware executing keyed GDSS.

\vspace{0.5em}
\noindent\textbf{Traffic-analysis residual.} Even with session-unique sync bursts, an adversary who knows the protocol family might search for energy events within the agreed epoch window. Keyed PN removes \emph{repeatable} correlation peaks across sessions (Figure~\ref{fig:cross}); it does not remove the existence of a physical transmission. Direction finding, TDOA, and power-change analytics remain in the ``physics-limited'' column of Table~\ref{tab:det}.

\vspace{0.5em}
\noindent\textbf{Post-quantum roadmap.} gr-linux-crypto exposes a hybrid KEM hook that currently raises \texttt{NotImplementedError} unless built with optional post-quantum support \cite{repo_crypto}. Migrating GR-K-GDSS session establishment from Brainpool ECDH to a PQ/classical hybrid would require new HKDF info labels, larger public keys on the wire (if sent), and updated latency budgets---none of which are exercised in the present repository state.

\vspace{0.5em}
\noindent\textbf{Regulatory and export context.} Covert waveforms may intersect spectrum-management and export-control rules depending on jurisdiction. This document contains no legal guidance; operators remain responsible for compliance, licensing, and dual-use classifications independent of the software licences of gr-k-gdss, gr-linux-crypto, and gr-qradiolink.

\section*{Appendix E: Security stack (from project README, abridged)}
\noindent\textbf{Layer 1 -- SOQPSK.} Continuous-phase modulation reduces discrete symbol transitions that would aid feature detectors.

\noindent\textbf{Layer 2 -- GDSS spreading/masking.} Spreading factor $N{=}256$ yields processing gain; keyed Gaussian masking aims for thermal-like PDF while requiring the session key to strip masks.

\noindent\textbf{Layer 3 -- LDPC.} Rates 1/2, 2/3, 3/4 with block lengths 576/1152/2304 bits (gr-qradiolink) recover part of the SNR penalty of GDSS relative to DSSS \cite{shakeel2023,chung2001}.

\noindent\textbf{Layer 4 -- ChaCha20-Poly1305.} RFC~8439 AEAD protects payload confidentiality and integrity before spreading \cite{rfc8439}.

\noindent\textbf{Layer 5 -- Brainpool ECDH.} BSI-influenced curve family; shared secret never sent on-air when public keys are authenticated out-of-band \cite{rfc5639,bsi2018}.

\noindent\textbf{Layer 6 -- Web of trust.} GnuPG signing chains bind public keys to verified identities; mitigates MITM on long-term keys.

\noindent\textbf{Layer 7 -- Kernel keyring.} Session subkeys may reside in the Linux keyring to limit user-space exposure; keys may clear when hardware tokens are removed \cite{repo_gdss}.

\noindent\textbf{Layer 8 -- Nitrokey HSM.} Private signing/decryption keys can remain on-device; PIN and brick policies address physical seizure \cite{repo_gdss}.

\noindent\textbf{Layer 9 -- Operations.} Mobility, short duty cycles, and traffic discipline raise the cost of long integrations against the noise floor \cite{repo_gdss}.

\vspace{0.5em}
\noindent\textbf{Appendix F: Fuzzing and formal methods (gr-linux-crypto).} The companion repository reports LibFuzzer campaigns exceeding 805 million executions without crashes, CBMC checks on core paths, and dudect-style timing experiments with $t$-statistics below published thresholds on representative MAC operations \cite{repo_crypto_results,reparaz2017}. These results validate the \emph{crypto library substrate}, not the radio-specific GDSS code path end-to-end.

\vspace{0.5em}
\noindent\textbf{Appendix G: IQ test artefacts.} The generator writes paired \texttt{.cf32} and metadata; analysis scripts check mean, variance symmetry, kurtosis, skewness, autocorrelation, KL divergence between standard and keyed transmissions, and cross-correlation of keyed sync bursts across synthetic sessions. All are software regressions; they do not replace calibrated laboratory spectrum measurements.

\section{Intellectual Property and Licensing}
The \textbf{GR-K-GDSS} (\texttt{gr-k-gdss}) software implementation is licensed under the \textbf{GNU General Public License v3.0 or later} (SPDX: \texttt{GPL-3.0-or-later}). C and C++ source files in that repository carry matching SPDX identifiers.

The project is \textbf{registered with the Open Invention Network} (OIN), a defensive patent pool that supports Linux-related ecosystems.

The \textbf{full text} of the GNU GPL version~3 is shipped in the GR-K-GDSS repository and can be read directly at:

\noindent\url{https://raw.githubusercontent.com/Supermagnum/GR-K-GDSS/main/LICENSE}

\noindent\textbf{Related repositories} (each documents its own license in-tree; comply with all applicable terms when building the stack):

\begin{itemize}\setlength{\itemsep}{2pt}
  \item \textbf{GR-K-GDSS:} \url{https://github.com/Supermagnum/GR-K-GDSS}
  \item \textbf{gr-linux-crypto:} \url{https://github.com/Supermagnum/gr-linux-crypto}
  \item \textbf{gr-qradiolink:} \url{https://github.com/Supermagnum/gr-qradiolink}
\end{itemize}

\begin{thebibliography}{99}
\bibitem{shakeel2023} I. Shakeel et al., ``Gaussian-Distributed Spread-Spectrum for Covert Communications,'' \textit{Sensors} 2023, 23(8), 4081. doi:10.3390/s23084081
\bibitem{rfc8439} Y. Nir, A. Langley, RFC 8439, ChaCha20-Poly1305, 2018.
\bibitem{rfc5869} H. Krawczyk, P. Eronen, RFC 5869, HKDF, 2010.
\bibitem{rfc5639} M. Lochter, J. Merkle, RFC 5639, Brainpool Curves, 2010.
\bibitem{bash2015} B. A. Bash \textit{et al.}, IEEE Commun. Mag., 2015.
\bibitem{bash2013} B. A. Bash \textit{et al.}, IEEE J. Sel. Areas Commun., 2013.
\bibitem{turner1991} L. F. Turner, NAECON, 1991.
\bibitem{gardner1992} W. A. Gardner, C. M. Spooner, IEEE Trans. Commun., 1992.
\bibitem{chung2001} S.-Y. Chung \textit{et al.}, IEEE Commun. Lett., 2001.
\bibitem{bernstein2008} D. J. Bernstein, SASC 2008.
\bibitem{bsi2018} BSI TR-03111, v2.10, 2018.
\bibitem{reparaz2017} O. Reparaz \textit{et al.}, CHES 2017 (dudect).
\bibitem{kaddoum2016} G. Kaddoum, IEEE Access, 2016.
\bibitem{itu1406} ITU-R P.1406, 2018.
\bibitem{stanag4539} STANAG 4539, HF channel models.
\bibitem{makhdoom2022} I. Makhdoom \textit{et al.}, Comput. Secur., 2022.
\bibitem{soltani2018} R. Soltani \textit{et al.}, IEEE Trans. Wirel. Commun., 2018.
\bibitem{repo_gdss} Supermagnum, ``gr-k-gdss,'' GitHub, 2025. \url{https://github.com/Supermagnum/GR-K-GDSS}
\bibitem{repo_crypto} Supermagnum, ``gr-linux-crypto,'' GitHub, 2025. \url{https://github.com/Supermagnum/gr-linux-crypto}
\bibitem{repo_crypto_results} gr-linux-crypto, \texttt{tests/TEST\_RESULTS.md}, 2026 snapshot.
\bibitem{repo_qrl} Supermagnum, ``gr-qradiolink,'' GitHub, 2025. \url{https://github.com/Supermagnum/gr-qradiolink}
\end{thebibliography}

\end{document}
